\clearpage
\begin{examplebox}
\textbf{例 1：不求两个根，直接求弦长。}
直线 $x+y-1=0$ 与椭圆 $x^2/8+y^2/4=1$ 的交点为 $P,Q$，求 $PQ$。
前面的投影图与直角三角形图即为本题模型。
\[
 \alpha=8,\ \beta=4,\ A=B=1,\ C=-1,\quad D=12,\ H=11>0.
\]
所以确有两个不同交点，通用公式给出
\[
 PQ=\frac{2\sqrt{8\cdot4\cdot(1^2+1^2)\cdot11}}{12}
   =\boxed{\frac{4\sqrt{11}}3}.
\]
\textbf{最低门槛的回验：} 代入 $y=1-x$ 得 $3x^2-4x-6=0$，
所以 $s=x_1+x_2=4/3,p=x_1x_2=-2$。于是
\[
 |x_2-x_1|=\sqrt{s^2-4p}=\frac{2\sqrt{22}}3,\qquad
 PQ=\sqrt2\,|x_2-x_1|=\frac{4\sqrt{11}}3.
\]
若另用未约分方程 $12x^2-16x-24=0$，其判别式为 $1408$；
约分方程的判别式为 $88$。分别除以二次项系数 $12$ 和 $3$ 后，根差相同。

\tcbbreak
\textbf{例 2：长轴换方向，求中点。}
直线 $2x+y-1=0$ 与椭圆 $x^2/3+y^2/4=1$ 交于 $P,Q$，求弦中点。
\begin{center}\begin{tikzpicture}[font=\small]
\draw[black!30,rounded corners=2pt] (0,0) rectangle (14.4,8.2);
\node[anchor=west,font=\small\bfseries] at (.5,7.72) {长轴在 $y$ 轴：仍按 $x^2,y^2$ 各自的分母代入};
% MODEL 3 4 2 1 -1
\begin{scope}[shift={(7.2,3.65)},scale=1.08]
\draw[blue,very thick] (0,0) ellipse (1.7320508076 and 2.0000000000);
\draw[-Stealth,black!45] (-2.3,0)--(2.3,0) node[right] {$x$};\draw[-Stealth,black!45] (0,-2.3)--(0,2.3) node[above] {$y$};
\coordinate (P) at (-0.4635254916,1.9270509831);
\coordinate (Q) at (1.2135254916,-1.4270509831);
\coordinate (M) at (0.3750000000,0.2500000000);
\coordinate (O) at (0.0000000000,0.0000000000);
\draw[orange!85!black,very thick] (P)--(Q);\fill[black] (P) circle (1.8pt) node[above left=6pt] {$P$};
\fill[black] (Q) circle (1.8pt) node[below right=6pt] {$Q$};
\fill[black] (M) circle (1.8pt) node[above right=7pt] {$M$};

\end{scope}
\node[anchor=west] at (.8,6.7) {$x^2/3+y^2/4=1,\quad 2x+y-1=0$};\node[anchor=west] at (.8,.6) {$M=(3/8,1/4)$：不求交点，也能得到弦中点};
\end{tikzpicture}\end{center}

直接按分母读出 $\alpha=3,\beta=4$，而不是按“较大者叫 $a$”换位。
$A=2,B=1,C=-1$，故 $D=16,H=15>0$。于是
\[
 M=\left(-\frac{2\cdot(-1)\cdot3}{16},-\frac{1\cdot(-1)\cdot4}{16}\right)
  =\boxed{(3/8,1/4)}.
\]
\textbf{回验：} 消去 $y$ 得 $16x^2-12x-9=0$，所以 $x_M=(12/16)/2=3/8$；
再由 $y_M=1-2x_M=1/4$。中点在直线上，但一般不在椭圆周上。

\tcbbreak
\textbf{例 3：先看几何，再用韦达交叉核验。}
直线 $y=kx+1$ 与椭圆 $x^2/2+y^2=1$ 交于不同两点 $P,Q$，
且 $OP\perp OQ$，$O$ 为原点，求 $k$。
\begin{center}\begin{tikzpicture}[font=\small]
\draw[black!30,rounded corners=2pt] (0,0) rectangle (14.4,8.2);
\node[anchor=west,font=\small\bfseries] at (.5,7.72) {定点已在椭圆上：先认出固定交点，几何短证更快};
% MODEL 2 1 1 1.4142135623730951 -1.4142135623730951
\begin{scope}[shift={(7.2,3.4)},scale=1.65]
\draw[blue,very thick] (0,0) ellipse (1.4142135624 and 1.0000000000);
\draw[-Stealth,black!45] (-1.9,0)--(1.9,0) node[right] {$x$};\draw[-Stealth,black!45] (0,-1.55)--(0,1.55) node[above] {$y$};
\coordinate (P) at (0.0000000000,1.0000000000);
\coordinate (Q) at (1.4142135624,0.0000000000);
\coordinate (S) at (-1.4142135624,0.0000000000);
\coordinate (O) at (0.0000000000,0.0000000000);
\draw[orange!85!black,very thick] (P)--(Q);\draw[purple,thick,dashed] (P)--(S);\draw[green!45!black,thick] (O)--(P);\draw[teal!80!black,thick] (S)--(Q);\pic[draw,angle radius=3mm] {right angle=P--O--Q};\fill[black] (P) circle (1.8pt) node[above right=6pt] {$P=(0,1)$};
\fill[black] (Q) circle (1.8pt) node[below right=6pt] {$Q_+$};
\fill[black] (S) circle (1.8pt) node[below left=6pt] {$Q_-$};
\fill[black] (O) circle (1.8pt) node[below left=7pt] {$O$};

\end{scope}
\node[anchor=west] at (.8,6.7) {$OP\perp OQ$ 迫使另一个交点落在 $x$ 轴上};\node[anchor=west] at (.8,.6) {$Q_\pm=(\pm\sqrt2,0)$，所以 $k=\pm1/\sqrt2$};
\end{tikzpicture}\end{center}

\textbf{优先短证：} 固定点 $(0,1)$ 已在椭圆上，可记为 $P$。
$OP$ 为竖直线，垂直条件迫使 $Q$ 在 $x$ 轴上。
椭圆上的横轴端点为 $Q_\pm=(\pm\sqrt2,0)$，所以
\[
 k=\frac{0-1}{\pm\sqrt2-0}=\boxed{\pm\frac1{\sqrt2}}.
\]
\textbf{韦达核验：} $A=k,B=-1,C=1,\alpha=2,\beta=1$，
$D=2k^2+1,H=2k^2$。$k=0$ 时相切，先排除；余下 $A\ne0,B\ne0$，可用两组公式：
\[
 x_1x_2=0,\qquad y_1y_2=\frac{1-2k^2}{2k^2+1}.
\]
$\overrightarrow{OP}\cdot\overrightarrow{OQ}=x_1x_2+y_1y_2=0$，得 $k^2=1/2$。
代回 $H=1>0$，两个交点确实不同，与短证和图形一致。
\end{examplebox}
